\section*{УВОД}
\label{sec:introduction}
\addcontentsline{toc}{section}{УВОД}

\subsection*{Актуалност}

Инжекцията е една от най-старите и най-устойчивите уязвимости в софтуера. Тя
продължава да заема видимо място в класациите на най-сериозните рискове за
уеб приложенията~\cite{owasp_top10} и е средство за документирани
киберпрестъпления, обсъдени в глава~\ref{sec:cybercrime}. Устойчивостта ѝ не се
дължи на липса на знание. Защитата срещу класическата SQL инжекция е известна
отдавна: стойностите трябва да се подават като параметри, а не да се вграждат в
текста на заявката~\cite{owasp_sqli_prevention}. Уязвимостта оцелява, защото
защитата зависи от дисциплината на всеки разработчик при всяка заявка.

Това е причината подходът към инжекциите да се измества от откриване и поправяне
на отделни случаи към тяхното изключване като клас. Съвместно предупреждение на
американските агенции CISA и FBI призовава производителите на софтуер да
елиминират SQL инжекциите именно по този начин~\cite{cisa_sqli_alert}. Възниква
въпросът какви средства позволяват такова изключване и докъде то стига.

\subsection*{Изследователски въпрос}

Настоящата работа изследва един такъв подход: конструиране на заявките по време
на компилация. При него структурата на заявката, имената на таблиците и на
колоните са фиксирани от типовата система, а стойностите могат да влязат само
през типизирани заместители. Работата отговаря на въпроса:

\begin{quote}
\textit{Може ли една абстракция, конструираща заявките по време на компилация, да
елиминира цяло семейство от класове инжекции при хетерогенни хранилища за
данни, и къде точно спира тази гаранция?}
\end{quote}

Въпросът е поставен за семейство от класове, а не само за SQL, защото
съвременните приложения използват различни видове хранилища, всяко със собствен
език и собствен клас инжекции. Втората част на въпроса е равностойна на първата.
Гаранция, чиито граници не са известни, не може да бъде използвана безопасно.

\subsection*{Цел и задачи}

Целта на работата е да провери, чрез измерване, а не само чрез разсъждение,
доколко конструирането на заявки по време на компилация предотвратява
инжекционните атаки при седем хетерогенни хранилища за данни, и да установи
границите на тази защита. За постигане на целта са поставени следните задачи:

\begin{enumerate}
    \item да се анализира инжекцията като средство за киберпрестъпления и
          правната рамка за противодействието ѝ;
    \item да се прегледат съществуващите защити и да се установи незапълнената
          нужда;
    \item да се формулира аргумент, свързващ формалното определение на
          инжекцията~\cite{su2006essence} с типовата система на C++;
    \item да се определи модел на заплахите и критерии за успех, фиксирани преди
          измерванията;
    \item да се анализира реализацията на гаранцията при всяко от седемте
          хранилища;
    \item да се изгради възпроизводим набор за оценка: отрицателни компилационни
          тестове, корпус от публикувани полезни товари и диференциален тест;
    \item да се оценят резултатите спрямо критериите и да се назоват границите на
          гаранцията.
\end{enumerate}

\subsection*{Обхват}

Обхватът е пътят на заявката от изграждането ѝ в приложението до предаването ѝ
на драйвера на хранилището. Извън обхвата са операциите по схемата, инжекцията от
втори ред, логическите грешки в приложенията, уязвимостите в самите драйвери и
отказът от услуга чрез скъпи заявки. Тези граници са обосновани в
точка~\ref{sec:tm-scope}.

Работата използва като платформа библиотека за обектно-релационно
картографиране, разработена в предходна дипломна работа за ОКС
„Бакалавър“~\cite{tsvetanov2026orm}. Предходната работа описва проектирането и
реализацията на библиотеката и не разглежда въпроса за сигурността. Настоящата
работа не претендира за нейните резултати. Тя изследва самостоятелен въпрос и
формулира приноси, които се отнасят само към него и към промените в
библиотеката, произтекли от изследването.

\subsection*{Структура на работата}

Глава~\ref{sec:cybercrime} разглежда инжекцията като средство за
киберпрестъпления и правната рамка за противодействието ѝ.
Глава~\ref{sec:existing} прави преглед на съществуващите защити и установява
незапълнената нужда. Глава~\ref{sec:theory} въвежда формалното определение на
инжекцията и аргумента за времето на компилация като граница на доверие.
Глава~\ref{sec:threat-model} определя модела на заплахите, критериите за успех и
метода на оценката. Глава~\ref{sec:enforcement} анализира как гаранцията се
постига при седемте хранилища. Глава~\ref{sec:case-study} проследява открит в хода
на работата пропуск в гаранцията, от откриването до поправката.
Глава~\ref{sec:evaluation} представя и тълкува резултатите от оценката.
Заключението обобщава приносите, ограниченията и насоките за бъдеща работа.
